\subsection{Partial operators}

In this section, K is a commutative field.

We recall (E, II, §3, p. 9–10) that a \(\text{graph}^{(2)}\) is a set all of whose elements are ordered pairs. If A and B are sets, a \(\text{correspondence}\) between A and B is a triple \((\Gamma, A, B)\), where \(\Gamma\) is a graph contained in A × B; its domain (also called \(\text{domain}\) ) is \(\text{pr}_1(\Gamma)\) , and the set of its values is \(\text{pr}_2(\Gamma)\) . A correspondence is a \(\text{function}\) (E, II, p. 13, def. 9) if its graph is functional and if its set of departure coincides with its domain of definition. Every subset of a functional graph is a functional graph.

DEFINITION 1. --- Let E and F be vector spaces over K. A partial operator u from E to F is a correspondence (Γ, E, F) between E and F satisfying the following conditions:

\begin{enumerate}
\def\labelenumi{(\roman{enumi})}
\item
  The graph \(\Gamma\) is a vector subspace of \(\mathrm{E} \times \mathrm{F}\) ;
\item
  The graph \(\Gamma\) is functional.
\end{enumerate}

If E = F, we say that u is a partial operator on E.

Let u be a partial operator from E to F. The graph \(\Gamma\) of the correspondence u is called the graph of the partial operator u, and is also denoted \(\Gamma_{u}\) . We denote \(\mathcal{P}(\mathrm{E};\mathrm{F})\) the set of partial operators from E to F; we simply denote \(\mathcal{P}(\mathrm{E}) = \mathcal{P}(\mathrm{E};\mathrm{E})\) .

Giving a partial operator from E to F amounts to giving a vector subspace D of E and a linear map u from D to F, the associated partial operator being the correspondence \((\Gamma, \mathrm{E}, \mathrm{F})\) where \(\Gamma \subset D \times F\) is the graph of u.

The domain of definition of a partial operator \(u\) is simply called the domain of \(u\) , and denoted \(\operatorname{dom}(u)\) .

Every linear map u from E to F is a partial operator from E to F.

If D ⊂ E is a vector subspace, we will denote by 1 \(_{D}\) the partial operator with domain D, which is the identity map on D, that is, the correspondence (Δ \(_{D}\) , E, E) where Δ \(_{D}\) is the diagonal of D × D (E, II, p. 13, def. 8). We shall denote by \(0_{\mathrm{D}}\) the partial operator with domain D that is zero on D, that is, the correspondence \((\mathrm{D} \times \{0\}, \mathrm{E}, \mathrm{F})\) .

Partial operators \(u = (\Gamma, \mathrm{E}, \mathrm{F})\) and \(u' = (\Gamma', \mathrm{E}, \mathrm{F})\) from E to F are equal if and only if \(\operatorname{dom}(u) = \operatorname{dom}(u')\) and if the linear maps \(u\) and \(u'\) from \(\operatorname{dom}(u)\) into F coincide.

Following E, II, §3, the following notions are defined:

\begin{enumerate}
\def\labelenumi{(\roman{enumi})}
\tightlist
\item
  The image of a subset A of E under u is the subset \(u(A \cap D)\) of F; it is simply denoted \(u(A)\) . The inverse image under u of a subset B of F is the subset \(\overline{u}(B)\) of D.
\end{enumerate}

If A (resp. B) is a vector subspace of E (resp. of F), then its image under u (resp. its inverse image) is a vector subspace of F (resp. of E).

The image of \(u\) is the vector subspace \(u(\mathrm{D})\) of \(\mathrm{F}\) , also denoted \(\operatorname{Im}(u)\) . One says that \(u\) is a surjective partial operator if \(\operatorname{Im}(u) = \mathrm{F}\) . The kernel of \(u\) is the vector subspace \(u^{-1}(\{0\})\) of \(\mathrm{E}\) , also denoted \(\operatorname{Ker}(u)\) . The kernel of \(u\) is reduced to 0 if and only if the linear map \(u\) of \(\operatorname{dom}(u)\) in \(\mathrm{F}\) is injective. We then say that \(u\) is injective. If \(u\) is injective and surjective, it is said to be bijective.

\begin{enumerate}
\def\labelenumi{(\roman{enumi})}
\setcounter{enumi}{1}
\item
  If E, F, and G are vector spaces over K and \(u = (\Gamma, \mathrm{E}, \mathrm{F})\) , \(v = (\Gamma', \mathrm{F}, \mathrm{G})\) are partial operators from E to F and from F to G, respectively, the composed correspondence \(v \circ u = (\Gamma' \circ \Gamma, \mathrm{E}, \mathrm{G})\) is a partial operator from E into G. Its domain is \(u^{-1}(\mathrm{dom}(v))\) . If H is a vector space over K and \(w = (\Gamma'', \mathrm{G}, \mathrm{H})\) a partial operator from G into H, we have \(w \circ (v \circ u) = (w \circ v) \circ u\) . One will sometimes write vu instead of \(v \circ u\) .
\item
  In particular, for every partial operator \(u\) from E into F and every \(a \in \mathrm{K}\) the partial operators \(au = (a1_{\mathrm{F}}) \circ u\) and \(ua = u \circ (a1_{\mathrm{E}})\) are defined. They are equal if \(a \neq 0\) or if the domain of \(u\) is equal to E; we have \(u0 = 0_{\mathrm{E}}\) and \(0u = 0_{\mathrm{dom}(u)}\) .
\end{enumerate}

Let E be a vector space. By the preceding, the set \(\mathcal{P}(\mathrm{E})\) endowed with the composition law defined by \((u,v)\mapsto u\circ v\) is an associative unital magma (A, I, p. 4, def. 5 and A, I, p. 12, def. 2) with identity element \(1_{\mathrm{E}}\) . For every \(n\in \mathbf{N}\) , we denote by \(u^n\) the composite \({}^n u\) (A, I, p. 13).

Moreover, we define the following notions:

\begin{enumerate}
\def\labelenumi{(\roman{enumi})}
\item
  If \(u = (\Gamma, \mathrm{E}, \mathrm{F})\) is a partial operator from E to F, and if G is a vector subspace of E, the reduction of \(u\) to G is the partial operator ( \(\Gamma \cap (\mathrm{G} \times \mathrm{F}), \mathrm{E}, \mathrm{F}\) ) from E to F. Its domain is \(\operatorname{dom}(u) \cap \mathrm{G}\) ; it will sometimes be denoted \(u|_{\mathrm{G}}\) when no confusion with the restriction of \(u\) to the subspace G is to be feared.
\item
  Let \(v\) an injective partial operator from F to E. The inverse correspondence \(v^{-1} = (\Gamma^{-1}, \mathrm{E}, \mathrm{F})\) of \(v\) is then a partial operator such that \(\operatorname{dom}(v^{-1}) = \operatorname{Im}(v)\) . One says that \(v^{-1}\) is the inverse partial operator of \(v\) . We have the equalities \(v \circ v^{-1} = 1_{\operatorname{dom}(v^{-1})}\) in \(\mathcal{P}(\mathrm{E})\) and \(v^{-1} \circ v = 1_{\operatorname{dom}(v)}\) in \(\mathcal{P}(\mathrm{F})\) . The partial operator \(v^{-1}\) is injective and we have \((v^{-1})^{-1} = v\) .
\item
  Let E, F, and G be vector spaces. Let u (resp. v) be an injective partial operator from E to F (resp. from F to G). Then the partial operator \(v \circ u\) is injective and \((v \circ u)^{-1} = u^{-1} \circ v^{-1}\) .
\item
  If u and v are partial operators from E to F, we say that v is an extension of u, and we write \(u \subset v\) if the graph of u is contained in the graph of v. This implies that \(\text{dom}(u) \subset \text{dom}(v)\) and that u is the reduction of v to \(\text{dom}(u)\) . The relation “ \(u \subset v\) ” is an order relation on \(\mathcal{P}(\text{E};\text{F})\) . We have for example \(au \subset ua\) for all \(a \in K\) and every \(u \in \mathcal{P}(\text{E};\text{F})\) .
\item
  Let E be a vector space over K and \((\mathrm{F}_{i})_{i\in\mathrm{I}}\) a family of vector spaces over K. For \(i\in I\) , let \(u_{i}\) a partial operator from E to \(F_{i}\) . The product partial operator of the \(u_{i}\) is the partial operator from E to the product vector space of the spaces \(F_{i}\) whose domain is the intersection D of the spaces \(\mathrm{dom}(u_{i})\) and which assigns to \(x\in D\) the family \((u_{i}(x))_{i\in\mathrm{I}}\) . It is denoted \((u_{i})_{i\in\mathrm{I}}\) .
\item
  Let A: F×F → F be the linear map \((x,y)\mapsto x+y\) . Let u and v be partial operators from E to F. The sum \(u+v\) is the partial operator A ∘ \((u,v)\) from E to F. Its domain is dom \((u)\cap\text{dom}(v)\) . For u, v, w in \(\mathcal{P}(\text{E};\text{F})\) , we have \((u+v)+w=u+(v+w)\) .
\end{enumerate}

Let G be a vector space over K. For all u and v in \(\mathcal{P}(\mathrm{E};\mathrm{F})\) and every \(w \in \mathcal{P}(\mathrm{F};\mathrm{G})\) , we have \(w \circ u + w \circ v \subset w \circ (u + v)\) . In general, there is no equality in this formula (exercise 1 of IV, p. 344), but this is the case when the domain of w is equal to F. For \(w \in \mathcal{P}(\mathrm{G};\mathrm{E})\) , we have \(u \circ w + v \circ w = (u + v) \circ w\) .

\begin{enumerate}
\def\labelenumi{(\roman{enumi})}
\setcounter{enumi}{6}
\tightlist
\item
  Let L be an extension of the field K. Let E and F be vector spaces over K and \(\mathrm{E}_{(\mathrm{L})} = \mathrm{L} \otimes_{\mathrm{K}} \mathrm{E}\) , \(\mathrm{F}_{(\mathrm{L})} = \mathrm{L} \otimes_{\mathrm{K}} \mathrm{F}\) the L-vector spaces obtained by extension of scalars from K to L (A, II, p. 82). For every partial operator u from E to F, we denote \(u_{(L)}\) the partial operator from \(\mathrm{E}_{(\mathrm{L})}\) into \(\mathrm{F}_{(\mathrm{L})}\) whose graph is the vector subspace \(L \otimes_{K} \Gamma_{u}\) of \(\mathrm{E}_{(\mathrm{L})} \times \mathrm{F}_{(\mathrm{L})}\) ; its domain is \(L \otimes_{K} \operatorname{dom}(u)\) and it coincides on this domain with the unique linear map that sends \(1 \otimes x\) onto \(1 \otimes u(x)\) for all \(x \in \operatorname{dom}(u)\) .
\end{enumerate}

Let v be a partial operator from E to F. We have \(u \subset v\) if and only if \(u_{(L)} \subset v_{(L)}\) .

\begin{enumerate}
\def\labelenumi{(\roman{enumi})}
\setcounter{enumi}{7}
\tightlist
\item
  Let \(E_{1}\) , \(F_{1}\) , \(E_{2}\) , \(F_{2}\) of K-vector spaces. Let u (resp. v) be a partial operator from \(E_{1}\) in \(F_{1}\) (resp. of \(E_{2}\) in \(F_{2}\) ). We denote by \(u \otimes v\) the partial operator from \(E_{1} \otimes F_{1}\) into \(E_{2} \otimes F_{2}\) with domain dom(u) \(\otimes\) dom(v) such that \((u \otimes v)(x \otimes y) = u(x) \otimes v(y)\) for all \((x, y) \in E_{1} \times E_{2}\) .
\end{enumerate}

